\documentclass[11pt]{article}
\usepackage[T1]{fontenc}
\usepackage{lmodern}
\usepackage[a4paper,left=25mm,right=30mm,top=22mm,bottom=24mm,marginparwidth=18mm]{geometry}
\usepackage{amsmath,amssymb,mathtools,stmaryrd}
\usepackage{tikz}
\usetikzlibrary{arrows.meta,calc,positioning}
\usepackage{caption}
\usepackage{xcolor}
\usepackage{marginnote}
\usepackage{microtype}
\usepackage{needspace}
\usepackage[hidelinks]{hyperref}
\definecolor{LeanTeal}{HTML}{087F72}
\definecolor{SourceGray}{HTML}{66736F}
\setlength{\parindent}{0pt}
\setlength{\parskip}{0.55em}
\setlength{\emergencystretch}{2em}
\newcommand{\leanmarker}[2]{%
  \Needspace{3\baselineskip}%
  \marginnote{\hyperlink{#1}{\color{LeanTeal}\fbox{\scriptsize Lean}}}%
  \hypertarget{#1}{\textbf{#2.}}%
}
\newcommand{\sourcepage}[1]{\par\medskip\noindent\textcolor{SourceGray}{\scriptsize Source PDF page #1}\par\smallskip}
\pagestyle{plain}
\begin{document}
\begin{center}\LARGE\bfseries Guide to the literature\par\medskip
\large Lean 4 Formalization Companion -- Chapter 9\end{center}
\medskip
\noindent\textit{This project-generated reading copy contains reviewed mathematical items, short source locators, and redrawn diagrams. It does not reproduce the original PDF or automatic OCR prose. Click a Lean tag in the web reader to open the mapped Lean declaration.}
\bigskip\hrule\bigskip
\section*{Section 9}
\sourcepage{152}
\leanmarker{def-9-1-1}{General reading guide} The source recommends systematic lecture notes by Scott (Sco81) and Plotkin (Plo81), the introductory exposition GS90, the continuous-lattice account GHK+80, Johnstone's treatment of Stone duality and continuous lattices (Joh82), and the introduction to partial orders DP90.

\leanmarker{def-9-1-2}{Reading for Section 2} Further reading: DP90 and Joh82.

\leanmarker{def-9-1-3}{Reading for Section 3} Further reading: Plo81, Gun92b, and Win93.

\leanmarker{def-9-1-4}{Reading for Section 4} Further reading: Jun89 and Jun90.

\leanmarker{def-9-1-5}{Reading for Section 5} Further reading: SP82, Fre91, Fre92, Pit93b, and Pit93a.

\leanmarker{def-9-1-6}{Reading for Section 6} Further reading: Plo76, Smy78, Win83, Hec91, and Sch93.

\leanmarker{def-9-1-7}{Reading for Section 7} Further reading: Abr90c, Abr91a, AO93, Ong93, Hen93, Bou94, Jen92, Jen91, and Smy83b.

\leanmarker{def-9-2-1}{Programming-language textbooks} The recommended semantics textbooks are Sch86, Ten91, Gun92b, and Win93.

\leanmarker{def-9-2-2}{Applications of Domain Theory} Additional applications are abstract interpretation and static analysis (Abr90a, BHA86, AJ91), databases (BDW88, BJO91), computational linguistics (PS84, PM90), artificial intelligence (RZ94), fractal image generation (Eda93b), and foundations of analysis (Eda93a).

\leanmarker{def-9-2-3}{Related Handbook chapters} The source lists six other Handbook chapters with substantial domain-theoretic content: Topology, Algebraic Semantics, Semantics of Types, Correspondence between Operational and Denotational Semantics, Abstract Interpretation, and Effective Structures.

\leanmarker{def-9-3-1}{Bibliography query API} Lean defines citation membership by
\[
  \operatorname{ContainsCitation}(G,k)\quad\Longleftrightarrow\quad
  \exists E\in G,\;\exists c\in E.\operatorname{citations},\;c.\operatorname{key}=k.
\]

\leanmarker{prop-9-3-2}{Verified guide queries} The executable bibliography API proves that Sco81 occurs in the complete reading guide and that Fre91 occurs in the Section 5 entry.

\end{document}
